% Part: many-valued-logic
% Chapter: infinite-valued-logics
% Section: introduction

\documentclass[../../../include/open-logic-section]{subfiles}

\begin{document}

\olfileid{mvl}{inf}{int}

\olsection{مقدمة}

ليس من شرط المصفوفة أن تكون قيم الصدق فيها متناهية. ومن أظهر
الخيارات لمجموعة لا متناهية مجموعة الأعداد النسبية المحصورة بين
$0$ و$1$، وهي~$V_\infty = [0,1] \cap \Rat$؛ أي
\begin{align*}
    V_\infty & = \Setabs{\frac{n}{m}}{n,m \in \Nat \text{ و } m>0 \text{ و } n\le m}.
\intertext{وعند النظر في مجموعة قيم الصدق غير المنتهية هذه، يفيد غالبًا
النظر أيضًا في المجموعات الجزئية}
V_m & = \Setabs{\frac{n}{m-1}}{n \in \Nat \text{ و } n<m}
\intertext{فمثلًا، $V_5$ هي المجموعة ذات قيم الصدق الخمس المتساوية التباعد،}
V_5 & = \{0, \frac{1}{4}, \frac{1}{2}, \frac{3}{4}, 1\}.
\end{align*}
وفي المنطقات المبنية على هذه المجموعات لا تكون، في العادة، إلا~$1$
مميزة؛ أي~$V^+ = \{1\}$. فـ~$1$ تقوم مقام الصدق المطلق، و$0$ مقام
الكذب المطلق، ويجوز مع ذلك لـ~!!{formula}s أن تأخذ أي قيمة متوسطة
في~$V$.

ويجوز بدل النسبية أن نأخذ~$V_{[0,1]} = [0,1]$، وهي مجموعة الأعداد
\emph{الحقيقية} بين~$0$ و$1$، أو أن نأخذ مجموعة جزئية أخرى لا
متناهية من~$[0,1]$. والمنطقات التي تتخذ مثل هذه المجموعات قيمًا
للصدق تسمى غالبًا \emph{ضبابية}.


\end{document}

